\chapter*{Abstract}

The proliferation of hyper-realistic deepfakes and AI-generated disinformation poses a severe threat to public trust and information integrity. This project presents \textbf{Satya}, a sophisticated multimodal verification assistant engineered to detect fake news and deceptive media. Unlike standard binary classifiers, Satya is designed as a Generative AI-powered system that analyzes text, extracts context from images, and inspects video for visual and audio manipulation cues. 

The core problem solved is the lack of reasoning in traditional detection mechanisms; users need to know \textit{why} a piece of news is fake, not just a probability score. To resolve this, our methodology employs a microservices-based architecture utilizing React and FastAPI, integrated with large language models (such as Google Gemini). The system performs input understanding, evidence retrieval, and generative verification, outputting a structured report encompassing a verdict, confidence score, detected claims, and supporting evidence. 

By leveraging GenAI for deep synthesis and asynchronous processing frameworks (RabbitMQ, Redis) for heavy video workloads, Satya represents a scalable, comprehensive approach to modern misinformation analysis. The resulting system successfully demonstrates the capacity to act as an AI-assisted verification tool, providing robust, explainable insights essential for countering digital deception in real-world scenarios.
